At the lowest point of its path, the star is on the opposite side of the celestial pole, toward the north.

The north celestial pole is at an altitude equal to the observer's latitude, \(\mathrm{lat}\). A star with declination \(\delta_*\) is \(90^\circ - \delta_*\) away from the north celestial pole. Therefore its lowest altitude is

\[
h_{\min} = \mathrm{lat} - (90^\circ - \delta_*)
\]

and therefore

\[
h_{\min} = \mathrm{lat} + \delta_* - 90^\circ
\]

The star is circumpolar when

\[
\mathrm{lat} + \delta_* > 90^\circ
\]

In that case its lowest point is still above the horizon. If this value is negative, the star is below the horizon at its lowest point.